%! Dividing Polynomials; Remainder and Factor Theorems — Unit 4, Lesson 4.3 — Exit Ticket
\documentclass[10pt]{article}
\usepackage{algebra2-article}
\usepackage{algebra2-boxes}
\newcommand{\TallMath}[1]{$\displaystyle #1\rule[-1.4em]{0pt}{3.2em}$}

\begin{document}

\pageheader{Unit 4, Lesson 4.3}{Exit Ticket}
\namedateperiod

\vspace{0.10in}

No notes. Watch your placeholders and your signs.

\begin{enumerate}[leftmargin=1.8em, itemsep=12pt]
  \item \textbf{Synthetic division, then finish.} Divide $f(x) = x^3 + 2x^2 - 5x - 6$ by $x + 1$.

    \vspace{3pt}
    The divisor is $x+1$, so $r = $ \blank{1.0cm}
    \begin{center}
    \renewcommand{\arraystretch}{1.6}
    \begin{tabular}{r|C{1.5cm}C{1.5cm}C{1.5cm}C{1.5cm}}
        &     &     &     &     \\
        &     &     &     &     \\ \cline{2-5}
        &     &     &     &     \\
    \end{tabular}
    \end{center}

    \vspace{2pt}
    Quotient: \blank{3.0cm} \quad Remainder: \blank{1.2cm}

    \vspace{3pt}
    Factor the quotient and write $f$ completely factored: \blank{5.0cm}

    \vspace{3pt}
    The three zeros of $f$ are \blank{4.0cm}

  \item \textbf{No dividing.} Let $g(x) = x^3 - 4x^2 + 2x + 5$.

    \vspace{3pt}
    What is the remainder when $g(x)$ is divided by $(x - 2)$? \blank{1.6cm}

    \vspace{3pt}
    Which theorem let you answer that without dividing? \blank{4.0cm}

    \vspace{3pt}
    Is $(x-2)$ a factor of $g$? \blank{1.2cm} \; How do you know? \blank{5.0cm}

  \item \textbf{SOL-style multiple choice.} Which binomial is a factor of
    $p(x) = x^3 - 6x^2 + 11x - 6$?
    \begin{enumerate}[label=\Alph*., leftmargin=1.9em, itemsep=1pt]
      \item $x + 1$
      \item $x - 1$
      \item $x - 6$
      \item $x + 6$
    \end{enumerate}
\end{enumerate}

\end{document}
